%! Counting and Probability with Counting — Unit 3, Lesson 3.5 — Homework (Answer Key)
% The lesson's SCORED (1 to 5) individual practice, the LAST component in the packet; assigned
% at Close & Assign and done outside class. DELTAMATH + JUSTIFICATIONS form: DeltaMath carries
% items 1-3 (plain \blank here, \ans in the key — an unboxed reference); students turn in items
% 4, 5, and 6, whose final answers are \gradedblank boxes copied verbatim to the key.
% THIRD CONTEXT: notes = student council at Fairview High; Guided Practice = a playlist; here,
% locker codes and a fantasy-team draft. TWO PAGES MAX: items 1-4 / items 5-6 and the spiralbox.
% Item 1 core FCP, repeats vs. no repeats. Items 2 and 3 the deliberate CONTRAST PAIR: the same
% n and r with roles and without. Item 4 the CRUX head-on, a justification. Item 5 probability
% with counting and an interpretation. Item 6 the SPIRAL (Lesson 3.4: conditional probability,
% dependent events), its boxed probabilities carrying the reason.
% THE NUMBERS (checked in Python): 10^4 = 10000; 10P4 = 5040; 15P3 = 2730; 15C3 = 455;
%   2730/455 = 6 = 3!; 6C3 = 20; 20/455 = 4/91 = 0.04396; P(2nd rookie | 1st rookie) = 5/14
%   vs P(1st rookie) = 6/15 = 0.4.
% PAGE LOCKSTEP: the key comes from templates/lesson/mkkey.py; work blocks and gradedblank lines
% are byte-identical; items 1-3 differ only by \blank -> \ans.
\documentclass[12pt]{article}
\usepackage{saar-article}
\usepackage{saar-key}   % pulls in -boxes; do NOT also load -boxes
\usepackage{grading}          % math-shared: \gradedblank, QR codes, marks
\gradingassignment{saar-u03-l05-hw}
\gradingcourse{stats_analysis_algebraic_reasoning}
\newcommand{\TallMath}[1]{$\displaystyle #1\rule[-1.4em]{0pt}{3.2em}$}
\newcommand{\Pnr}[2]{{}_{#1}P_{#2}}
\newcommand{\Cnr}[2]{{}_{#1}C_{#2}}
% Word bank strip. Defined per document; the key inherits it and prints the same strip.
\newcommand{\wordbank}[1]{\par\vspace{2pt}\noindent\fcolorbox{goldacc}{goldbg}{\parbox{\dimexpr\linewidth-2\fboxsep-2\fboxrule\relax}{\small\textbf{Word bank:} #1}}\par\vspace{4pt}}
% Handwriting room inside every \begin{work} block; moves the blank and key together.
\setlength{\workrowsep}{8pt}

\begin{document}

\pageheader{Unit 3, Lesson 3.5}{Homework --- Answer Key}
{\small\itshape Write each final answer in its box. Only what is inside a box is graded.}\par
% NAMESTRIP: no \namedateperiod here — mirrors the blank. NO teachernote here —
% teacher prose lives in the lesson plan.

\vspace{0.08in}

% Authored identically in homework/ and homework_key/.
\begin{remindbox}
\small
\textbf{This is your graded homework.} It is scored 1 to 5 and due at
the start of the first class after two study halls. It is assigned today and done on your own,
outside class --- flip back to the \emph{Keep in Mind} box on the cover and to rows 2--4 of
your notes.
\end{remindbox}

\vspace{0.06in}

\boxguard[20]
\begin{scenariobox}[Locker Codes and a Fantasy-Team Draft]{cerulean}
School lockers open with a $4$-digit code, each digit $0$ through $9$. In a class
fantasy-sports league, each team builds its roster from the same pool of $15$ players;
$6$ of the $15$ are rookies. Show your work outside the boxes.
\end{scenariobox}

\vspace{0.06in}

\boxguard
\begin{notesbox}{Counting Codes and Rosters}
\wordbank{permutation \;\textbullet\; combination}
\begin{enumerate}[leftmargin=1.6em, itemsep=6pt, topsep=3pt]

  \item How many $4$-digit codes are possible \textbf{if digits may repeat}?
        \hfill \ans{10,000}

        \begin{work}
          10 \cdot 10 \cdot 10 \cdot 10 &= 10{,}000 \text{ codes}
        \end{work}

        How many \textbf{if no digit may repeat}?
        \hfill \ans{5,040}

        \begin{work}
          10 \cdot 9 \cdot 8 \cdot 7 &= 5{,}040 \text{ codes}
        \end{work}

  \item Your team fills $3$ ranked spots from the pool of $15$: \textbf{starter, backup, and
        reserve}. Name the count, then find it.
        \par\vspace{4pt}
        Count: \ans{permutation} \hfill Number of ways: \ans{2,730}

        \begin{work}
          \Pnr{15}{3} &= \dfrac{15!}{12!} \\
                      &= 15 \cdot 14 \cdot 13 = 2{,}730
        \end{work}

  \item Now you pick $3$ of the same $15$ players for your bench. \textbf{All bench spots are
        alike.} Name the count, then find it.
        \par\vspace{4pt}
        Count: \ans{combination} \hfill Number of ways: \ans{455}

        \begin{work}
          \Cnr{15}{3} &= \dfrac{15 \cdot 14 \cdot 13}{3 \cdot 2 \cdot 1} \\
                      &= 455
        \end{work}

\end{enumerate}
\end{notesbox}

\boxguard[30]
\begin{notesbox}{Counting Codes and Rosters, continued}
\begin{enumerate}[leftmargin=1.6em, itemsep=5pt, topsep=2pt, start=4]

  \item Devon says, \emph{``Items 2 and 3 both pick $3$ of $15$, so the answers must be the
        same.''} Say what Devon did not check, and explain why item~2's answer is exactly $6$
        times item~3's.
        \gradedresponse{\gradedblank[id=q4, type=rubric, points=2, display={Devon did not check whether order matters. Swapping two ranked players changes the roster, but swapping two bench players does not. Each group of 3 can fill the 3 ranked spots in 3! = 6 ways, so 2,730 = 6 x 455.}]{\linewidth}{2in}}

  \item A random draw gives a team $3$ of the $15$ players. Find the probability that \textbf{all
        $3$ are rookies}, as a decimal to three places.
        \hfill \blankbox{\gradedblank[id=q5a, type=short, points=1, answer={0.044}, check=numeric, tolerance=0.001, display={0.044}]{2.2cm}{18pt}}

        \begin{work}
          P &= \dfrac{\Cnr{6}{3}}{\Cnr{15}{3}} = \dfrac{20}{455} = \dfrac{4}{91} \approx 0.044
        \end{work}

        If $91$ teams were drawn this way, predict how many would be all rookies, and say what
        that means about an all-rookie team.
        \gradedresponse{\gradedblank[id=q5b, type=rubric, points=1, display={About 4 of the 91 teams would be all rookies. An all-rookie team is unusual, about 4 teams in every 91, but it can happen.}]{\linewidth}{2in}}

  \item \emph{Back to Lesson 3.4.} Two players are drawn from the $15$ one at a time, without
        replacement. Decide whether the events are independent or dependent, and let the two
        probabilities below be your reason. Write each probability as a fraction.
        \par\vspace{6pt}
        $P(\text{first is a rookie})$: \blankbox{\gradedblank[id=q6a, type=short, points=0.5, answer={6/15}, check=expression, display={6/15}]{1.9cm}{22pt}} \hfill $P(\text{second is a rookie} \mid \text{first is a rookie})$: \blankbox{\gradedblank[id=q6b, type=short, points=0.5, answer={5/14}, check=expression, display={5/14}]{1.9cm}{22pt}}
        \par\vspace{8pt}
        The two events are: \blankbox{\gradedblank[id=q6c, type=short, points=1, answer={dependent}, check=text]{3.2cm}{18pt}}

\end{enumerate}
\end{notesbox}

\boxguard[5]   % the short spiralbox needs about five baselines; a larger guard risks a third page
\begin{spiralbox}
\textbf{Next class: Unit 3 Review.} This was the last lesson of Unit~3. Next class we review for
the Unit~3 test: two-way tables, probability models and rules, conditional probability, and
today's counting. Item~4 is exactly the kind of question the test asks.
\end{spiralbox}

\end{document}
